\chapter{Evaluating a Family}\label{ch:evaluation}

A variant family can fail in more ways than a single film. Its optics can leak, its realizations can be incoherent alone or together, its switching can strand viewers, its audience can feel that it attended several events rather than one, and its equipment can do things its wearers never agreed to. This chapter sets out how a family and its exhibition should be evaluated, in four groups of tests: technical, narrative, social, and ethical. Appendix~D collects them as a protocol.

\section{Technical evaluation}

The optical system is measured first, because everything else depends on it.

\emph{Timing.} The display's refresh must be stable, the glasses must hold phase against the synchronization signal, and drift must be measured over a full screening, not a few minutes. Phase error that is harmless at the start of a film can accumulate.

\emph{Leakage.} Crosstalk is measured in both of the senses \Cref{ch:phase} distinguished. The optical measure is the fraction of light from other streams reaching the eye. The informational measure is whether leaked content is legible: whether test viewers can identify faces, words, or objects flashed in other streams at the family's actual leakage level. The informational test is the one that matters for protected variation, and it must be run at the moments the protection profile marks as sensitive, with the blanking and allocation settings used there.

\emph{Luminance and flicker.} Each stream's luminance is measured through the glasses, in coincident and divergent passages, and flicker is assessed centrally and peripherally at those luminances. Brightness changes from adaptive gating and allocation shifts are logged, and their timing is checked to confirm they fall at cuts or ramp gradually, and never recur as a regular pulse.

\emph{Seating.} All of these are measured across the auditorium, not only at its center: viewing angle affects some shutter technologies, and distance and screen size affect flicker and legibility.

\emph{Failure.} Every exhibition must declare what the glasses do when they lose the timing signal or their power. There are three choices. Glasses that go transparent show the composite; if the composite is designed, this is a safe failure, and if the family contains protected variation, it may expose it. Glasses that go opaque interrupt the film but reveal nothing. Glasses that fall back to a declared common stream keep the viewer in the film but may place them in a realization they did not choose. The right choice depends on the family's protection profile and may change during the film. What is not acceptable is an undeclared failure mode, discovered by an audience during a screening.

\section{Narrative evaluation}

Each realization is watched alone and judged as a complete film. The variant matrix checks of \Cref{ch:matrix} establish coherence; this evaluation establishes quality, which no check can.

The family is then evaluated as a family. Reviewers confirm that its realizations stand in the intended relations: that the genre pair tells one story, that the depth family's explanations serve their viewers, that the endings are each earned. They look for unintended asymmetry, such as one realization given the better writing or the stronger performances, which is especially important when one realization is an access stream.

Switching is tested directly. Reviewers switch at the points the switch system allows, to confirm that the moves work as intended, and at points it forbids, to see what goes wrong and whether the prohibition is worth what it costs the viewer. They try paths with several switches, since \Cref{ch:switch-admissibility} showed these can fail when single switches do not.

Finally, reviewers watch different realizations in one room and compare their memories afterward. The comparison shows what each realization actually established for its viewers, which may differ from what the matrix says it establishes, for the reason \Cref{ch:depth} gave: presentation is not uptake.

\section{Social evaluation}

A family can be technically clean and narratively strong and still fail at what the book has argued is its purpose: to be one occasion experienced differently. That failure can only be found by watching audiences.

The definition of an occasion-preserving family in \Cref{ch:synchronized-difference} gives measurable proxies. Bounded drift can be measured from the editor's data. The frequency of structural rendezvous can be read from the switch system. The presence of a common channel is a fact about the exhibition. But the definition's success is a claim about experience, and it must be tested by asking audiences.

Relevant questions include whether viewers felt they had attended the same event as the others in the room; whether hearing others' reactions to a different realization felt like connection, confusion, or intrusion; whether the conversation afterward was enriched or frustrated by the difference; whether viewers felt ranked or sorted by their choice; whether companions coordinated their choices, and why; and, for families with a designed composite, whether anyone who saw it understood it as part of the work. Preference ratings alone will not answer these questions. A family can be enjoyed by every viewer and still have failed to be one occasion, if each experienced it as a private film projected on a public wall.

\section{Ethical audit}

The last group of tests concerns the commitments of Part~VIII, and it is the easiest to omit, because nothing in the film breaks if it is skipped.

The glasses are inspected. The claim of \Cref{ch:phase}, that the simplest glasses can be shown to record and measure nothing, is only as good as the inspection that confirms it for the glasses actually handed out: no cameras, no eye trackers, no storage beyond the dial state, no transmitter beyond what timing requires.

The registry and manifest of \Cref{ch:persuasion} are checked against what was exhibited, stream by stream.

The provenance of every realization is confirmed: which were made by the work's authors, which were added or altered afterward, and whether the alterations are declared.

The disclosure given to the audience is checked for what it actually says: that variation exists, what the realizations are, how each viewer's realization is chosen, and whether it can be changed.

A family that passes the first three groups of tests and fails this one should not be exhibited. It may be excellent cinema. It is also, in the terms of the previous chapters, an instrument for telling different people different things without their knowledge, and the quality of the cinema makes that worse, not better.
